% s-230-p230-m003; original PDF 230
\Aphorism{s-230-p230-m003}{325}
Ohne eigentlich künstlerisch zu sein, hat unser Zeitalter für die Kunst doch viel Zeit und Ernst übrig; vielleicht rührt das daher, dass das Ideal Wissenschaft nachgerade nicht nur land-, sondern beinahe gassenläufig geworden ist und sich mit jeder Sorte Mensch einlässt. Ehedem, als noch Religion die Rolle dieser Aphrodite Pandemos spielte, war umgekehrt die Wissenschaft eine Zuflucht derer, die einsam und vornehm leben wollten. Es wäre möglich, dass der Kreis sich schlösse und Übersättigung an der Kunst wieder zu einer Art Religion hinführte; statt Culturgeschichte dürfte man dann genauer sagen: Flucht des Einzelnen vor der nachdrängenden Masse.
\par

